\chapter{Conclusion and Limitations}
The project collected \Papers{} official main-track AAAI and IJCAI papers from
2021--2025 and constructed a graph of \TotalNodes{} nodes and
\TotalEdges{} edges. DOI-based enrichment, position-specific researcher
identification and paper-specific affiliations supply the bibliographic base.
A frozen MiniLM classifier adds multi-label topics with abstention; accepted
paper evidence supports collaboration and topic-profile relationships.

The analysis reveals an expanding selected corpus, a large coauthorship component
alongside many small groups, repeated academia--industry publication links,
and different organization topic portfolios. Community and centrality views
help inspect collaboration structure. Citation analysis identifies Informer as
the most cited paper within this restricted corpus. The graph and Explorer
provide a practical way to move from such summaries to the supporting entities
and papers.

The main limitations are the two-venue boundary, 11 unmatched works, withheld
researcher positions, incomplete affiliations, seven unresolved historical DOI
discrepancies, inferred-topic uncertainty and citations restricted to the corpus.
These affect visible connectivity, portfolio denominators and rankings.
Publication collaboration records historical co-publication, not current
employment or contractual relationships. Centrality describes graph structure.

The most useful next steps are targeted person/source checks, improved
paper-specific affiliation coverage, and a fresh topic reference set before
stronger topic interpretations. Broader venue coverage and age-aware citation
comparisons would extend the analysis. Any data correction should preserve
the original evidence and document the affected records.

Full technical evidence and reproduction instructions are linked in the project overview.
